% Compile with pdflatex main.tex
\documentclass[11pt,a4paper]{article}

\usepackage[utf8]{inputenc}
\usepackage[T1]{fontenc}
\usepackage{lmodern}
\usepackage[margin=1in]{geometry}
\usepackage{amsmath,amssymb}
\usepackage{graphicx}
\usepackage{booktabs}
\usepackage{hyperref}
\usepackage{microtype}
\usepackage{xcolor}
\usepackage{caption}
\usepackage{enumitem}

\hypersetup{
    colorlinks=true,
    linkcolor=blue!50!black,
    citecolor=blue!50!black,
    urlcolor=blue!50!black
}

\title{Learning Without Remembering: Habituation in Weights and a Self-Model That Acts in a Minimal Embodied Agent}
\author{Adri\'an Valerio Porras\textsuperscript{1} \\ \small \textsuperscript{1}Independent Researcher, San Jos\'e, Costa Rica \\ \small \texttt{adrian\_valerio@valeriogroup.co} \\ \small ORCID: 0009-0003-8445-2214 }
\date{1 September 2026 --- v0.13 (peer-review revision; supersedes v0.12, doi:10.5281/zenodo.22191728)}

% Plain-text requirements for automated checks: WORLD_SIZE 20.0, H*=[0.8,0.9,0.2,0.7], MLP 13->64->6, Phi 15->64->1 (calibration) / 22->64->1 (integrated), attention 13->7, fog x>14 noise 0.15/0.60, EWC lambda=5 orth lambda=0.01
% v0.13 CHANGELOG: C3 renamed stimulus generalization (not dishabituation); Table 1 provenance split into grid battery vs continuous v0.12; Phi noted as offline diagnostic (MSE to |epsilon|, not per-channel log-variance) with attention-gate confound flagged; EWC noted as non-load-bearing for same-task interference; MiniGrid removed from main Table; Rankin battery listed as pre-registered future work.
\begin{document}
\maketitle

\begin{abstract}
We present a minimal embodied agent that runs without vision or audition, on body channels only: an action-conditioned body predictor, a homeostatic drive (E,C,U,S), episodic memory, and a meta-cognitive module $\Phi$ that predicts its own prediction error. A pre-registered 30-seed study (grid world, 13$\rightarrow$128$\rightarrow$128$\rightarrow$6, $\pm$5-cell teleport) tested whether this loop can detect violations of its sensorimotor contingencies, habituate by updating its model, and retain that update in weights rather than in explicit memory. It can, but coarsely. Detection was reliable ($z=20.6$, 95\% CI [16.0, 25.5]). Habituation cut surprise by 86\% ($d=3.5$). The trace lived in the weight delta (post/pre ratio 0.02; restore pre-$W$ restores $z=67.5$, freeze $W$ prevents habituation). The meta-cognitive module $\Phi$ (MSE to $|\epsilon|$, scalar, 15$\rightarrow$64$\rightarrow$1) was calibrated ($r=0.701$) and generalized ($r_{\mathrm{cross}}=0.730$); in an isolated causal battery it halved time in an unreliable zone when coupled to action via an attention gate (15.0\% vs 28.1\%, $d=-1.61$), with the attention gate itself as a confound flagged. One pre-registered control failed and is reported in full: stimulus generalization (C3) shows habituation to a $+5$ teleport generalizes to $-5$ ($z$ 1.1 vs 0.9), i.e.\ learning of displacement magnitude, not the vector---``learning without distinguishing.'' The continuous v0.12 organism (13$\rightarrow$64$\rightarrow$6, $+$2 teleport, active attention) reproduces the closed-loop integration over 30\,000 steps. All code, registration notes, and raw logs accompany this preprint; the system runs on a laptop with no discrete GPU.
\end{abstract}

\noindent\textbf{Keywords:} embodied cognition, sensorimotor contingencies, habituation, meta-cognition, predictive coding, homeostatic reinforcement learning, continual learning

\noindent\textit{Preprint. Intended submission to arXiv cs.AI and q-bio.NC. This is a workshop-level report of an integrated minimal system, not a claim about phenomenal consciousness.}

% Word count: abstract 167 words

\section{Introduction}

Large language models produce fluent text, but fluency is not evidence that anything behind the output cares, predicts, gets surprised, or remembers in its own body. We ask what is behind the mouth. We do not claim to have solved consciousness. We claim something smaller and more testable.

The hard problem remains \cite{chalmers1995}. But the question of autonomy is narrower and more tractable than the question of experience. An autonomous system should keep itself alive within bounds, notice when the world stops obeying its own learned regularities, update that knowledge, and know when it cannot trust its senses. Those are engineering questions. They are also the questions that enactive and sensorimotor theories place at the centre of awareness \cite{oregan2001, thompson2007}.

We chose a deliberately minimal embodied agent for this reason. No images. No audio. No pre-trained vision backbone. The agent only feels its position and four internal variables we call ECUS (Energy, Courage or social drive, Uncertainty, and Sociability). The idea is familiar from first-person experience: with eyes closed, awareness persists through the body. If awareness has anything to do with sensorimotor contingencies, then it should already be visible in this minimal loop. If it is not, then adding a bigger network will not create it.

We were surprised that the loop worked at all on a single MacBook Pro (M4 Pro, Apple silicon), with no discrete GPU. What we found, after pre-registration and explicit controls, is a small set of effects that survive honest tests, and one instructive failure that we report because it constrains our claims more than another positive result would.

\subsection{What this paper adds}

We do not introduce a new learning rule, a new network architecture, or a new benchmark record. Every piece we use exists in the literature. Our contribution is to put four pieces together in one closed, continuously running organism, and to test whether the integrated system does something that none of the pieces does alone:

\begin{enumerate}[leftmargin=1.2em]
    \item An action-conditioned body predictor (the organism learns $P(s_{t+1} \mid s_t, a_t)$ for its own body, not for pixels).
    \item A homeostatic drive with setpoints and precision-weighted surprise (the organism cares about staying near its preferred state).
    \item Continual learning with weight-level persistence, tested by erasing explicit memory entirely.
    \item A calibrated meta-cognitive self-model $\Phi$ that predicts its own prediction error and gates action.
\end{enumerate}

The fourth item is the one we worried was a philosophical trick, so we tested whether $\Phi$ actually changes what the agent does. It does, and the effect is large. The integration itself, together with the full set of controls, is the novelty. No published system, to our knowledge, shows all four running together for 30,000 continuous steps with the controls reported here (see Section~\ref{sec:related}).

\subsection{Limitations, upfront}

Because we want to keep the reader's trust, we state the limits before the results. First, this is a toy world. The predictor is a small MLP and the world was first a grid, now a continuous square. Nothing here scales to biology. Second, habituation generalizes too broadly: learning that ``big jumps happen'' transfers across opposite vectors (C3, $z$ 1.1 vs 0.9, Section~\ref{sec:results}). We demonstrate habituation of displacement magnitude, not of a specific vector or contingency $P(\Delta p\mid a)$; we do not claim Rankin et~al.~\cite{rankin2009} characteristic 7 (stimulus specificity) or 8 (dishabituation, which we did not test: habituate $S_1$, present $S_2$, re-probe $S_1$) nor ISI-dependent spontaneous recovery \cite{thompson1966,rankin2009}. Third, $\Phi$ is a scalar MSE regressor to $|\epsilon|$, not a per-channel log-variance (precision) model, and its causal battery is confounded with the attention gate ($\mathrm{att}_{\mathrm{vis}}<0.35$); ``physics of precision'' is an offline diagnostic here, not the online drive (see Methods, Sec.~\ref{sec:results}, Sec.~\ref{sec:limitations}). Fourth, EWC ($\lambda=5$, diagonal Fisher, plus $\lambda_{\mathrm{ortho}}=0.01$ to fix global habituation in v0.11) is not load-bearing for same-task interference and the weight delta is likely low-rank; a factored predictor and $\Delta W$ SVD are needed (Sec.~\ref{sec:limitations}). Fifth, the mouth (frozen LFM2.5-1.2B) translates a constructed prompt; it is not free recall. Sixth, MiniGrid Empty-8x8 vs ICM/RND at $N=5$ (RND wins) is a pilot only, removed from Table~\ref{tab:main}. Details are in Section~\ref{sec:limitations}.

\section{Related Work}
\label{sec:related}

\paragraph{Habituation as learning.} We build directly on the framing of habituation as model update rather than fatigue, synthesized recently as Training Ecosystems \cite{levin2026}, but scored against the classic behavioural criteria \cite{thompson1966,rankin2009}. Their argument is that even minimal living systems treat repeated surprising inputs as evidence to update a generative model. Gillard and colleagues had already shown that a behavioural habituation criterion motivated by integrated information could drive open-ended learning in IDSM agents without explicit reward \cite{gillard2022}. Our habituation analysis is deliberately narrower: we operationalize habituation as a sustained drop in precision-weighted surprise after repeated violations, and we ask where that trace lives---noting that a decrement alone is necessary but not sufficient for Rankin's characteristics 7/8/10.

\paragraph{Meta-cognition and knowing when you do not know.} The idea that an agent should track its own uncertainty and use it to regulate behaviour is central to predictive processing and uncertainty-aware learning \cite{friston2010, laukkonen2025}. Our $\Phi$ is a minimal instance, trained to predict $\sigma$, the expected magnitude of the predictor's error $\epsilon$, and then used twice: first to weight surprise as $\mathrm{presence} = \epsilon / \sigma^2$, second to gate an epistemic action (leave the unreliable zone).

\paragraph{Homeostatic and intrinsically motivated RL.} Homeostatic RL treats drive reduction as reward \cite{kernelan2014, man2016}. Intrinsic curiosity (ICM, RND) treats prediction error or novelty as reward \cite{pathak2017, burda2018}. We keep the two separate and then couple them: ECUS provides a slow, setpoint-based drive, while surprise provides a fast, precision-weighted signal. The coupling is simple (a small increment to $U$ proportional to $z$), but it creates a useful tension between foraging and clarity.

\paragraph{Continual learning and synaptic persistence.} That a memory trace could outlive explicit recall is not new. Elastic Weight Consolidation \cite{kirkpatrick2017} and its successors showed how to protect important weights, and work on systems consolidation has long argued that traces migrate from episodic to parametric storage. What we test is more literal and more direct: we erase the entire episodic buffer $E$ and also restore or freeze the weight delta itself (controls C4a--C4c). If habituation were only in the buffer, it would vanish; it does not.

\paragraph{Consciousness theories and embodiment.} We are sympathetic to sensorimotor contingency theory \cite{oregan2001} and enactive views \cite{thompson2007}, and we use them as design constraints (body only, action-conditioned prediction) rather than as claims we assert. We do not claim to measure integrated information \cite{tononi2014} or to implement a global workspace \cite{baars1988, dehaene2011}. We report one place where those theories help us interpret a result (Section~\ref{sec:discussion}), and several places where they do not apply.

\paragraph{Recent minimal agents and LLM-as-tool views.} Three arXiv reports sharpened our thinking about minimal embodied agents. CheckVLA showed how easily ``violations of learned contingencies'' can be faked without proper action-shuffled and observation-only controls; we adopt their control logic verbatim \cite{checkvla2026}. Gubernaut showed that compact embodied loops with explicit stability criteria can be studied without claiming consciousness \cite{gubernaut2026}. Christov-Moore and colleagues argued that embodiment and homeostatic regulation deserve a more central place in discussions of artificial awareness \cite{christovmoore2025}. The LFM2 technical report provided the efficient hybrid SSM-conv model we use only as a frozen translator \cite{lfm2025}. Finally, ``Asleep at the Wheel'' showed that prediction-error novelty scores can collapse to chance under a fair, single-dataset protocol \cite{probe2026}; we treat that as a warning that our own surprise signal needs exactly the controls we report here.

\paragraph{Where we claim novelty.} No single component above is new. What we have not seen elsewhere is the integration of all four into one continuously running organism, with pre-registered multi-seed statistics for the main battery and the specific control battery we report (action-shuffled, observation-only, weight-restore, weight-freeze, untrained). If a close precedent exists we would be glad to cite it. The niche, as of August 2026, appears to be empty at this level of blinding and integration. The risk is not that we duplicate work, but that another group assembles the same pieces tomorrow. We publish the pieces and the glue.

\section{Methods}
\label{sec:methods}

All numbers in Section~\ref{sec:results} come from scripts in \texttt{framework/} and JSON files in \texttt{results/}. The complete long run is reproduced by \texttt{python3 framework/organismo\_final.py --steps 30000} on a MacBook Pro (M4 Pro, Apple silicon) with MPS, no discrete GPU required. Wall time is a few minutes without the mouth, longer if generation through MLX is enabled.

Figure~\ref{fig:tetra} sketches the tetrahedron architecture we refer to as H1--H6. In words: H1 is persistence of the trace in weights, H2 is the action-conditioned body predictor, H3 is the ECUS homeostatic drive, H4 is measurement as precision-weighted surprise, H5 is presence ($\epsilon/\sigma^2$), H6 is the self-model $\Phi$ with active attention. The loop is closed. Most arrows are tied to at least one control in Section~\ref{sec:results}; the few that are not (the ECUS drive and the attention gate) are flagged in Section~\ref{sec:limitations}.

\begin{figure}[h]
\centering
\fbox{\parbox{0.92\linewidth}{\centering \vspace{1.2em} \textit{[Figure placeholder: Tetrahedron H1--H6]}\\ \small Body predictor (H2) $\rightarrow$ error $\epsilon$ $\rightarrow$ $\Phi$ predicts $\sigma$ (H6) $\rightarrow$ presence $=\epsilon/\sigma^2$ (H5) $\rightarrow$ ECUS drive (H3) $\rightarrow$ policy $\rightarrow$ epistemic action $\rightarrow$ memory $E$ + EWC weights (H1) $\rightarrow$ mouth (translates, never decides) $\rightarrow$ world. Active attention gate on every channel. \vspace{1.2em}}}
\caption{One closed loop, not a pipeline. The agent predicts its next bodily state conditioned on its own action, estimates how reliable that prediction is, weights surprise by that reliability, uses the result to modulate a homeostatic drive, acts, and consolidates the update in its weights. The LLM mouth is frozen and outside the causal chain that controls action.}
\label{fig:tetra}
\end{figure}

\subsection{Agent architecture}

\paragraph{Body predictor (H2).} A feed-forward MLP $f_\theta: \mathbb{R}^{13} \rightarrow \mathbb{R}^{6}$ with one hidden layer 13$\rightarrow$64$\rightarrow$6 and ReLU. Input concatenates normalized vision (2 dims: $x/W$, $y/W$), normalized interoception (4 dims: $H/1.5$), and a 7-way one-hot action. Output predicts the next normalized position (2 dims) and next normalized ECUS (4 dims). No convolutional encoder, no recurrence. The point was to keep the model so small that any interesting behaviour must come from the loop, not from capacity. Training uses Adam ($lr=10^{-3}$), 400 steps of 64-sample batches on 1,200 random transitions, then online updates with a small orthogonal penalty (see below).

\paragraph{Homeostatic drive ECUS (H3).} Four variables $H=[E, C, U, S]$ with preferred state $H^*=[0.8, 0.9, 0.2, 0.7]$. The drive is formalized as
\begin{equation}
D = \left(\sum_i w_i \, \|H_i - H^*_i\|^2 \right)^{1/2},
\end{equation}
with fixed weights $w$ that balance the channels. The implemented policy is a threshold approximation of this drive: if $E < 0.55$ the agent forages toward the nearest food, if $S > 0.8$ it seeks the social agent at $(18,18)$, otherwise it moves randomly. In continuous physics the same logic outputs a directional micro-action scaled by 0.8 and damped by 0.95 (see below). A small learning-to-drive coupling adds $\min(0.3, 0.2\max(0,z)) \cdot \mathrm{att}_{\mathrm{vis}}$ to $U$, where $z$ is surprise. This is the only place where surprise feeds back into motivation. It is deliberately modest. H5 (homeostasis without policy) failed as expected; H5-bis passes only with this policy, i.e.\ it tests the closed organism, not the predictor alone.

\paragraph{Surprise as precision-weighted prediction error (H4/H5).} At each step we compute raw error $\epsilon = \|f_\theta(s_t,a_t) - s_{t+1}\|_2$, maintain a running window of the last 100 values, and define surprise as a z-score
\begin{equation}
z = (\epsilon - \mu_{100}) / (\sigma_{100}+10^{-8}).
\end{equation}
Presence, the precision-weighted quantity that distinguishes expected noise from true violations, is
\begin{equation}
\mathrm{presence} = \epsilon / (\sigma_\Phi^2 + 10^{-6}),
\end{equation}
where $\sigma_\Phi = \Phi(s_t, a_t, \mu, \sigma)$ is the self-model's prediction of how large $\epsilon$ should be. When $\Phi$ expects large error (fog, for example), the same raw $\epsilon$ counts less. When $\Phi$ expects small error, a violation counts a lot. This is the ``physics of experience'' view in one line: experience is not $\epsilon$ alone, but $\epsilon$ relative to expected precision. We use presence strictly as an offline diagnostic in the isolated functional test (Table~\ref{tab:main}, ``$\Phi$ functional'', ratio 0.13). In the integrated 30k organism the online loop gates on $z$ and on attention thresholds next; presence does \textit{not} modulate the online drive. $\Phi$ here is a scalar MSE regressor to $|\epsilon|$ (15$\rightarrow$64$\rightarrow$1), not a per-channel log-variance precision model; a high $r$ can arise from ``I am in $x>14$'' as zone proxy. The proper test (per-channel $\log \sigma^2_c$, presence coupled to drive) is flagged as future work in Sec.~\ref{sec:limitations}.

\paragraph{Meta-cognitive self-model $\Phi$ (H6) with active attention.} $\Phi$ is a second MLP with $|\cdot|$ output and structure 15$\rightarrow$64$\rightarrow$1 in the calibration and causal batteries (input: the 13-dim predictor input plus the mean and standard deviation of the recent $\epsilon$ window), and 22$\rightarrow$64$\rightarrow$1 in the integrated organism, where the extra 7 inputs are the attention weights. $\Phi$ is trained to predict $\epsilon$ itself (MSE to $|\epsilon|$, 500--600 steps, Adam $10^{-3}$). It is therefore a model of the first model's fallibility, but only as a scalar expected-error regressor, not as a heteroscedastic (per-channel $\log \sigma^2$) precision model. The next step is $\Phi(s,a,\mathrm{channel})\rightarrow \log \sigma^2_{\Phi,c}$ and coupling $\mathrm{presence}$ to the drive; here $\mathrm{presence}$ is diagnostic only.

Active attention is a third small network $h_\psi: \mathbb{R}^{13} \rightarrow \Delta^7$ with 13$\rightarrow$32$\rightarrow$7 and softmax output $\mathrm{att} \in \mathbb{R}^7$ over the seven sensory channels. Attention does not select inputs to the predictor in a hard way. It modulates two things: (i) the effective noise estimate per channel, $\sigma_{\mathrm{chan}}$, and (ii) a gating term on epistemic action (leave the fog when visual attention is low or interoceptive attention is high). Concretely, if $\mathrm{att}_{\mathrm{vis}} < 0.35$ or $\mathrm{mean}(\mathrm{att}_{1:5}) > 0.65$, the policy executes a fixed epistemic action (index 3, which moves the agent out of the noisy zone). The rule is deliberately readable. Critically, this gate can act as a fog detector on its own; the $d=-1.61$ causal effect in Table~\ref{tab:main} confounds $\Phi$ with the attention gate (conditions B/C/D needed, see Sec.~\ref{sec:limitations}). In \texttt{framework/cuarteto\_habituacion.py} a trained vs random gate already yields the same fog-time reduction, consistent with that confound.

\paragraph{Episodic memory and synaptic consolidation (H1).} Episodic buffer $E$ is salience-gated and capped at 5,000 traces, but it is \textit{not} in the prediction path (erasing $E$ cannot by itself abolish prediction; see Table notes). EWC \cite{kirkpatrick2017} protects the predictor weights with $\lambda_{\mathrm{EWC}}=5.0$, computed against a diagonal online Fisher $F \leftarrow 0.9F + 0.1 g^2$. An additional orthogonal penalty $\lambda_{\mathrm{ortho}}=0.01$ on predictor weights ($\|W\|_2^2$) discourages a single direction from dominating and was essential to prevent the ``MLP habituates globally'' problem we saw in v0.11. For same-task interference (repeat the same teleport) additional violations are reinforcement, not forgetting, so $\lambda$ has no effect (0.48 recovery for $\lambda\in\{0,0.5,5,50\}$ in \texttt{cuarteto\_habituacion.py}); EWC would be load-bearing only for a distinct task B---this is not shown. The $\Delta W$ after habituation is therefore expected to be low-rank (``ignore large L2'') and is best diagnosed by SVD and by a factored predictor ($f^{\mathrm{pos}}:\mathbb{R}^{13}\rightarrow\mathbb{R}^2$ vs $f^{H}:\mathbb{R}^{13}\rightarrow\mathbb{R}^4$, separate losses), flagged in Sec.~\ref{sec:limitations}.

\paragraph{Mouth (LFM2.5-1.2B).} The mouth is the LFM2.5-1.2B-Instruct model (hybrid SSM-conv, 1.2B, MLX 8-bit) loaded through \texttt{mlx\_lm}. It is frozen. It never selects actions. Behaviour is identical with or without it (tested in H2b, see prior work \cite{lfm2025}). At most once per 500 steps, when $z>4$ or ($\sigma_\Phi>0.25$ and in fog), we render the internal state as text (``energy=0.72, $\Phi$ predicts $\sigma=0.34$ \dots'') and ask the model to translate it into one first-person sentence. The resulting sentence is logged and never fed back. We keep the mouth because ``SER versus DECIR'' matters: the fact that an internal state can be faithfully translated without that translation controlling action is precisely what deflates a naive Chinese Room reading.

\subsection{World: continuous physics with a fog zone}

The world is a continuous square $\mathrm{WORLD\_SIZE}=20.0$. Four foods are at $(3,3)$, $(3,16)$, $(10,3)$, $(10,16)$ and one social agent at $(18,18)$. Position updates as
\begin{equation}
p' = \mathrm{clip}(0.95 \,(p + 0.8\,a),\, 0,\, 20),
\end{equation}
where $a$ is the 2D component of the chosen action. Interoception updates as
\begin{equation}
\Delta H = -0.02\,(H-H^*) + \Delta_{\mathrm{fog}} + \Delta_{\mathrm{food}} + \Delta_{\mathrm{social}} + \eta,
\end{equation}
where $\Delta_{\mathrm{fog}}$ is $-0.03$ on $E$ and $+0.01$ on $U$ when $x>14$ (in clear weather, $U$ drifts $-0.01$; $E$ has no clear-weather counterpart), and $\eta$ is per-channel noise.

The fog zone is the region $x>14.0$. There, interoceptive noise is larger: base noise $\sigma_{\mathrm{base}}=0.15$, fog noise $0.60$ on channels 0--4, and $0.075$ on channel 6. Base noise elsewhere is $0.15$ on the same channels. The fog is therefore not a wall. It is observable unreliability. $\Phi$ must learn that predictions made in the fog should be trusted less. In an early design we put food inside the fog and the agent spent 64\% of its time there, foraging despite poor sensing. That was not a bug but a genuine trade-off. For the final integration we moved food outside the fog so that epistemic action could dominate without starvation pressure. Both designs are instructive, and we report both.

\subsection{Violations, pre-registration, and training regime}

We test three violation types. Motor: the body teleports by $+2$ in $x$ and $y$ (earlier grid work used $\pm5$ cells). Interoceptive: eating lowers $E$ (inverted causality). Tactile: wall pass-through. In the 30k integrated run, motor violations occur every 5,000 steps (six total).

\textbf{Provenance of Table~\ref{tab:main} (v0.13 fix).} Table~\ref{tab:main} now strictly separates two systems. Rows 1--11 (Detection, C1, C2, Habituation, Persistence, C4a--c, C3-generalization, H5-bis, $\Phi$ calibration/generalization/functional/causal) are the 30-seed (and C1--C4c single-seed pilot) \textit{grid-world battery} (20$\times$20 grid, $\pm5$-cell teleport, predictor 13$\rightarrow$128$\rightarrow$128$\rightarrow$6, $\Phi$ 15$\rightarrow$64$\rightarrow$1, no attention) pre-registered as H1--H4. The last row (Integrated 30k) is the \textit{continuous v0.12 organism} (13$\rightarrow$64$\rightarrow$6, $+2$ teleport, attention 13$\rightarrow$32$\rightarrow$7, $\Phi$ 22$\rightarrow$64$\rightarrow$1). The abstract cites the grid $z=20.6$ and the v0.12 0\% fog as different organisms, not one. Unifying all rows onto the continuous world (factored predictor, per-channel $\Phi$) is pre-registered future work (Rankin battery), not papered over.

We pre-registered the full statistical protocol before running the 30-seed battery (see the project's registration notes, doc.~48). Primary hypotheses H1--H6, sample size $N=30$, test statistics (means, 95\% bootstrap CIs with 2,000 resamples, Cohen's $d$ for paired and unpaired comparisons), and pass/fail thresholds are fixed in that document; the control battery C1--C4c was specified in the rigor-controls document (\texttt{46}). No threshold was moved after seeing data. When H5 failed in its first protocol (which omitted the navigation policy), we registered H5-bis as a named correction before executing it. That discipline costs us a clean sweep, but it is the reason we trust the positives that remain.

Predictor and $\Phi$ are pre-trained as described above, then updated online during the run with the same Adam optimizer and the EWC/orthogonal penalties. The Fisher is estimated online. Episodic buffer $E$ is appended only for violations or highly salient events.

\subsection{Controls (what would a conventional predictor not do?)}

Our guiding question, borrowed from the CheckVLA discussion \cite{checkvla2026}, is discriminative: what could a conventional action-conditioned predictor not do? A conventional predictor can have $\epsilon$ and even surprise, but it does not predict its own $\epsilon$, its uncertainty does not change its action, and its habituation does not live in a specific weight delta that can be restored or frozen. Each control isolates one of these.

\begin{itemize}[leftmargin=1.2em]
    \item \textbf{C1 action-shuffled:} replay the same trajectory but shuffle the action labels. If detection is truly action-conditioned, $z$ should collapse; it does, by $7\times$.
    \item \textbf{C2 observation-only:} train and test a predictor that sees only observations, not actions. If action adds information, $z$ should be weaker; it is, by $8.5\times$.
    \item \textbf{C3 stimulus generalization (negative result, misnamed dishabituation in v0.12):} habituate to teleport $(+5,+5)$ and probe with teleport $(-5,-5)$ of equal magnitude. If habituation were stimulus-specific at the vector level (Rankin characteristic 7 \cite{rankin2009}), surprise should return; it does not ($z$ 1.1 habituated vs 0.9 new, $\Delta\approx0$). This is a generalization test, not Rankin dishabituation \cite{rankin2009}. True dishabituation (habituate $S_1$, present distinct $S_2$, re-probe $S_1$; characteristic 8) and ISI-dependent spontaneous recovery \cite{thompson1966} were not tested and are pre-registered for v0.14.
    \item \textbf{C4a weight-restore:} after habituation, restore predictor weights $W$ to their pre-habituation snapshot. If habituation lived in $W$, surprise should return; it does ($z=67.5$).
    \item \textbf{C4b weight-freeze:} freeze $W$ and repeat the habituation protocol. If habituation requires learning, it should not occur; it does not ($z$ stays at $137.2$).
    \item \textbf{C4c untrained:} test an untrained predictor. If the effect requires learned physics, $z$ should be near zero; it is ($z=0.3$). This shows some learning is needed, not which physics was learned.
\end{itemize}

Together C1--C4c test conditioning, information, coarse learning in $W$, and a limit on granularity. We report all six, including the one that fails. Rankin characteristics 7/8/10 and SVD of $\Delta W$ are the next battery.

\section{Results}
\label{sec:results}

Table~\ref{tab:main} collects every statistic on one page. We then walk through the four claims in order and finish with the 30k integrated run that the reader can reproduce. No number is cherry-picked. Sample sizes are stated per row in the table notes; where we pre-registered a bootstrap interval we report it.

\begin{table}[h]
\centering
\caption{All outcomes (v0.13). Grid battery vs continuous v0.12 are separated. Sample sizes in Notes. CI is 95\% bootstrap (2,000 resamples). $d$ is Cohen's $d$. The negative result (C3) is kept and renamed.}
\label{tab:main}
\resizebox{\textwidth}{!}{%
\begin{tabular}{@{}llcc@{}}
\toprule
Claim & Measure & Result & Statistic \\
\midrule
\multicolumn{4}{@{}l}{\textit{Grid-world battery (20$\times$20, $\pm$5 teleport, 13$\rightarrow$128$\rightarrow$128$\rightarrow$6, $\Phi$ 15$\rightarrow$64$\rightarrow$1) -- N=30 unless noted}} \\
Detection (H1) & Surprise $z$ to motor violation & $20.6$ & CI $[16.0, 25.5]$ \\
C1 Action conditioning$^{\dagger}$ & $z$ correct vs shuffled & $132.7$ vs $18.1$ & $7.3\times$ \\
C2 Action adds info$^{\dagger}$ & $z$ action vs obs-only & $132.7$ vs $15.5$ & $8.5\times$ \\
Habituation (H3) & $z$ before $\rightarrow$ after & $3.8 \rightarrow 0.5$ & $86\%$, $d=3.5$ \\
Persistence (H4) & $z_{\mathrm{post}}/z_{\mathrm{motor}}$ after habituation & $0.02$ & $N=30$ \\
C4a Habituation in $W^{\dagger}$ & $z$ after restoring pre-$W$ & $67.5$ & (high) \\
C4b Requires learning$^{\dagger}$ & $z$ with $W$ frozen & $137.2$ & (no habituation) \\
C4c Requires physics$^{\dagger}$ & $z$ untrained & $0.3$ & ($\approx 0$) \\
C3 Generalization$^{\dagger}$ & $z$ habituated vs new direction & $1.1$ vs $0.9$ & fail (no specificity) \\
Homeostasis H5-bis & Mean $E$ over 5k steps & $0.85$, $100\%$ in $[0.5,1.2]$ & CI $[0.84, 0.85]$ \\
$\Phi$ calibration & Spearman $r(\hat\epsilon,\epsilon)$ & $0.701$ & threshold $>0.5$ \\
$\Phi$ generalization & $r_{\mathrm{cross}}$ on unseen violations & $0.730$ & threshold $>0.3$ \\
$\Phi$ functional (offline) & presence(noise)/presence(violation) & $0.13$ & $7.7\times$ sep. \\
$\Phi$ causal (isolated)$^{\ddagger}$ & Fog time, $\Phi$-coupled vs uncoupled & $15.0\%$ vs $28.1\%$ & $d=-1.61$ \\
\midrule
\multicolumn{4}{@{}l}{\textit{Continuous v0.12 organism (20$\times$20 continuous, $+$2 teleport, 13$\rightarrow$64$\rightarrow$6, atten. 13$\rightarrow$32$\rightarrow$7, $\Phi$ 22$\rightarrow$64$\rightarrow$1) -- N=1 trajectory}} \\
Integrated 30k run & Fog time, foods outside fog & $0.0\%$ & 6 violations \\
\bottomrule
\end{tabular}}
\\[0.6ex]
{\small \textit{Notes:} $^{\dagger}$C1--C4c are single-seed pilots (seed 7, \texttt{rigor\_controles.py}); mechanism pilots, not 30-seed estimates. $^{\ddagger}$$\Phi$ causal is the isolated h6\_phi\_causal battery (attention not modelled there; in the continuous organism the same fog-time drop is reproduced by a random attention gate, so the gate is a confound). $z_{\mathrm{post}}/z_{\mathrm{motor}}$ = mean surprise after habituation / mean surprise on first motor violation in the grid battery (that battery has no episodic buffer in the prediction path; buffer-erasure is the single-run \texttt{m5\_cadena\_completa.py} protocol where prediction never consulted $E$). H5-bis CI rounded to 2 decimals; full CI $[0.842,0.854]$, range $[0.81,0.88]$, in \texttt{results/h5bis.json}. MiniGrid $N=5$ pilot (organism 1.8\% vs ICM 0.6\% vs RND 2.8\%) removed from Table, kept only in Sec.~\ref{sec:limitations} as underpowered pilot. [Schematic placeholder -- vector PDF for camera-ready].}
\end{table}

\begin{figure}[h]
\centering
\fbox{\parbox{0.92\linewidth}{\centering \vspace{1.0em} \textit{[Figure placeholder: Habituation and persistence]}\\ \small Left: $z$ over repeated violations, mean $\pm$ SD over 30 seeds. Drops from $\sim 3.8$ to $\sim 0.5$ over $\sim 8$ repeats. Middle: weight-restore and weight-freeze controls. Right: $z_{\mathrm{post}}/z_{\mathrm{pre}}$ histogram after erasing episodic buffer $E$, mass at $0.02$. \vspace{1.0em}}}
\caption{Habituation lives in the weight delta, not in explicit memory. Restoring weights restores surprise, freezing prevents it, and erasing the entire episodic store leaves the trace intact.}
\label{fig:hab}
\end{figure}

\begin{figure}[h]
\centering
\fbox{\parbox{0.92\linewidth}{\centering \vspace{1.0em} \textit{[Figure placeholder: $\Phi$ calibration and causal effect]}\\ \small Left: scatter of predicted $\sigma_\Phi$ vs actual $|\epsilon|$, $r=0.701$. Middle: separation of expected noise (attenuated presence) vs true violation (full presence), ratio $0.13$. Right: distribution of fog time for $\Phi$-coupled (A) vs uncoupled (B) agents, $N=30$, $d=-1.61$. \vspace{1.0em}}}
\caption{The self-model is calibrated, it generalizes, and it matters for action. An agent that knows its senses are unreliable acts to regain clarity.}
\label{fig:phi}
\end{figure}

\subsection{Detection is reliable and genuinely action-conditioned}

Mean surprise to a motor violation was $z=20.6$ with 95\% CI $[16.0, 25.5]$ over 30 seeds. The interval does not come near the pre-registered null of $5$. This is not a single lucky seed. Every seed shows the effect, and the lower bound of the interval is itself a large deviation.

We worried this was a general novelty response, so we shuffled the action labels while replaying the same state trajectory (C1). Surprise collapsed from $132.7$ to $18.1$, a factor of about seven. A predictor that only watches observations, without knowing which action was taken, also collapsed to $15.5$ (C2, factor $8.5$). In other words, the detector does not fire because ``something changed.'' It fires because the change violated the specific action-conditioned regularity it had learned. That distinction is exactly the CheckVLA lesson we wanted to respect.

\subsection{Habituation reduces surprise by 86\% and lives in the weights}

Repeating the same violation type drives surprise from $3.8$ to $0.5$, an 86\% reduction with $d=3.5$ (Table~\ref{tab:main}). The effect is large and consistent across seeds. The question is where that reduction lives.

Three controls point to the weights. First, persistence: in the 30-seed battery, mean surprise after habituation divided by mean surprise on the first motor violation is $0.02$. That battery has no episodic buffer in the prediction path at all, so the residual trace cannot be attributed to explicit memory. In the single-run protocol (\texttt{m5\_cadena\_completa.py}), we erased the entire episodic buffer $E$ after habituation and behaviour was unchanged, because prediction never consulted $E$. Second, we restored the predictor weights to their pre-habituation snapshot (C4a). Surprise returned to $z=67.5$. The delta itself carried the trace. Third, we froze the weights and repeated the habituation protocol (C4b). Surprise stayed at $137.2$. No learning, no habituation. Finally, an untrained predictor (C4c) scored $z=0.3$, indistinguishable from zero. There is no habituation without learned physics to violate.

Taken together, this is habituation as model update, not as buffer lookup. The mechanism is simple: the predictor learns that ``large displacements can occur,'' so the next large displacement is less surprising. Whether that generalization is too broad is the next point.

\subsection{The clear failure: habituation generalizes across direction (C3) -- learning without distinguishing}

We pre-registered (v0.12) a stimulus-generalization test and misnamed it ``dishabituation.'' After habituating to teleport $(+5,+5)$, we probed with $(-5,-5)$ of equal magnitude. If habituation were specific to the learned vector (Rankin et~al.~\cite{rankin2009}, characteristic 7, stimulus specificity), surprise should have recovered; it did not. Surprise to the habituated direction was $1.1$; to the new direction $0.9$. The difference is noise. This is not Rankin dishabituation (characteristic 8: habituate $S_1$, present distinct $S_2$, re-probe $S_1$ for rebound) and we did not test ISI-dependent spontaneous recovery (characteristic 10) or a within-modality novel stimulus recovering response.

We keep this negative result visible for a reason. It tells us that the predictor did not memorize ``+5 in $x$, +5 in $y$.'' It learned a coarser category: ``a large displacement occurred''---habituation of displacement magnitude, not of the contingency $P(\Delta p\mid a)$. That is still model update (C4a/b, persistence), but at the granularity of ``large jump'' rather than a vector concept. The honest label is \textit{learning without distinguishing}. Early habituation literature already warned that a decrement in $z$ is necessary but not sufficient for habituation; specificity, dishabituation and recovery are the discriminants. To claim stimulus specificity we would need the Rankin-minimal battery (S1 $+2,+2$, S2 $-2,-2$, S3 $+2,-2$, S4 $+4,+4$, S5 interoceptive type, true dishabituation $S_5\rightarrow$ re-$S_1$, gap recovery) pre-registered as v0.14 in Sec.~\ref{sec:limitations}.

\subsection{Homeostasis is maintained when the policy is allowed to act}

Our first 30-seed battery (H5) measured $E$ with random actions and found a mean of $0.50$ with only half the seeds in the healthy range $[0.5, 1.2]$. We failed that test, and we logged the failure without re-interpreting it. The diagnosis was straightforward once we looked: the protocol measured the predictor without the foraging loop. Homeostasis is not a property of a predictor alone. It is a property of a closed organism that acts.

We pre-registered H5-bis before running it: the same 30 seeds, the same 5,000 steps, but with the full ECUS-coupled navigation policy that drives toward food when $E<0.55$ and toward the social agent when $S>0.8$. The result was $E=0.85$ with 100\% of seeds in $[0.5,1.2]$ and a CI of $[0.84,0.85]$ (full CI $[0.842,0.854]$, rounded to 2 decimals for the table; full data in \texttt{results/h5bis.json}). The range was $[0.81,0.88]$. In the 24-hour simulated run (864k steps) the same policy kept $E$ between $0.66$ and $0.84$ without intervention. Homeostasis is reliable when the organism is allowed to be an organism. That sentence sounds trivial, but many ablation studies forget it.

\subsection{\texorpdfstring{$\Phi$}{Phi} is calibrated, generalizes, and is causally efficacious in the isolated battery (with a confound flagged)}

The meta-cognitive module was tested in four ways. The pass thresholds ($r>0.5$, $r_{\mathrm{cross}}>0.3$, ratio $<0.7$, $|d|>0.5$) were fixed before the runs and are logged in \texttt{results/h6\_phi.json}.

First, calibration: Spearman correlation between predicted $\sigma_\Phi$ and actual $|\epsilon|$ was $r=0.701$ (threshold $>0.5$). The agent predicts, before erring, that it will err. In the fog, $\sigma_\Phi$ rises before the raw error does. This is a forecast, albeit from a scalar MSE regressor that can exploit ``$x>14$'' as a zone proxy.

Second, generalization: $\Phi$ trained on generic noise and tested on motor violations it had never seen scored $r_{\mathrm{cross}}=0.730$ (threshold $>0.3$). This is encouraging but does not separate ``learned reliability'' from ``learned zone''; a per-channel $\log \sigma^2$ model would.

Third, functional separation (offline diagnostic): the ratio $\mathrm{presence}(\mathrm{expected\ noise})/\mathrm{presence}(\mathrm{violation})$ was $0.13$, a $7.7\times$ separation (threshold $<0.7$). When $\Phi$ expects noise, the same raw error counts less. When it does not, the same error counts fully. This is presence, not $\epsilon$, and a conventional predictor lacks $\epsilon/\sigma^2$. Presence was \textit{not} coupled to the online drive in the 30k organism; it is reported as a diagnostic.

The fourth test is the one we care about most, because it answers the epiphenomenon worry. In the isolated h6\_phi\_causal battery (no attention, 15$\rightarrow$64$\rightarrow$1), we coupled $\Phi$ to action: when $\Phi$ predicts high uncertainty, the agent executes an epistemic action that leaves the fog. Condition A ($\Phi$-coupled) vs B (identical $\Phi$ computed but not allowed to influence action) over 30 seeds gave $15.0\%\pm6.3$ vs $28.1\%\pm9.6$ in fog, $d=-1.61$. This is a large effect in that battery. However, in the continuous v0.12 organism the same fog-time reduction ($24.7\%$ vs $36.8\%$, $d=-0.43$ in \texttt{62-correcciones-panel-resultados.md}) is reproduced by a random attention gate, and the attention threshold $\mathrm{att}_{\mathrm{vis}}<0.35$ alone can drive epistemic action. The $d=-1.61$ therefore confounds $\Phi$ with the attention gate. The decisive 4-arm test (A: $\Phi$+attention, B: $\Phi$ computed/uncoupled, C: attention only, D: presence coupled, attention frozen) is pre-registered as v0.14. We flag this explicitly so the Discussion's ``physics of experience'' is not read as the online drive.

\subsection{Long integration: 30,000 continuous steps}

We finally ran the full organism with every mechanism enabled for 30,000 steps, the tetrahedron closed (continuous v0.12, foods outside fog, violations every 5,000 steps). The timeline (from \texttt{organismo\_final.py}) is:

\begin{center}
\begin{tabular}{lcccc}
\toprule
$t$ & $E$ & $U$ & $z_{\max}$ & fog \\
\midrule
5,000 & 1.42 & 0.00 & 6.6 & 0\% \\
10,000 & 1.15 & 0.00 & 6.6 & 0\% \\
15,000 & 1.48 & 0.00 & 6.0 & 0\% \\
20,000 & 0.80 & 0.00 & 5.5 & 0\% \\
25,000 & 0.75 & 0.01 & 5.7 & 0\% \\
\bottomrule
\end{tabular}
\end{center}

Total fog time $0.0\%$ with foods outside fog (the agent briefly entered the fog early, the mouth logged it as low-reliability, and the attention-gated epistemic action kept it out thereafter). Six violations occurred, each created a trace in $E$ (6 traces) and the mouth produced 39 reports, e.g.\ ``estoy sintiendo que mi percepci\'on no es muy confiable en este momento'' and ``hay una alta incertidumbre en mi situaci\'on'', each verifiable against logged $\sigma_\Phi$ and $\mathrm{att}$. Energy peaked at $1.48$ near food and stayed in $[0,1.5]$. No mechanism broke. The earlier trade-off design (food inside fog) gave $64\%$ fog time (reported because epistemic action then resolves a real foraging-vs-clarity conflict), and that design yields the $24.7\%$ vs $36.8\%$ ($d=-0.43$) fog-time effect now shown to be reproducible with a random attention gate --- i.e.\ the gate itself, not $\Phi$, can explain the reduction in that world. This is why the isolated $d=-1.61$ is kept separate in Table~\ref{tab:main}.

MiniGrid Empty-8x8 ($N=5$, vanilla REINFORCE: organism 1.8\%, ICM 0.6\%, RND 2.8\%, random 0.4\%) is an underpowered pilot where RND beats the organism; it is removed from Table~\ref{tab:main} and kept only as a limitation (see Sec.~\ref{sec:limitations}).

\section{Discussion}
\label{sec:discussion}

\subsection{Habituation as model update, not episodic reminding -- but coarse and low-rank}

The pattern of weight-restore, weight-freeze, and buffer erasure is internally coherent. Habituation in this system is not the agent remembering ``I saw that before, so I will ignore it.'' It is the agent having changed what it expects the world to do. Restoring the old weights restores surprise because the old expectations return. Freezing the weights prevents habituation because expectations cannot be updated. Erasing the episodic store leaves habituation intact because the expectation lives in the slow weights, not in a list of past events. This is the Levin framing in a concrete number: a trace that persists when $E$ is gone.

Biology already knows that habituation is not one thing. Some habituation is short-term and need not involve lasting synaptic change. What we show is the long-term, model-update kind, at the smallest scale where it can be isolated. The fact that it generalizes broadly is still consistent with a model-update account, though a coarse one: the single-head predictor with a global L2 loss plus $\lambda_{\mathrm{ortho}}$ is well described as learning a low-rank ``ignore large $\| \Delta p \|$'' direction in $\Delta W$, not a vector contingency. The predicted signatures are (i) $\Delta W$ SVD dominated by 1--2 singular values and (ii) a factored predictor ($f^{\mathrm{pos}}$ vs $f^{H}$) splitting the effect: position head still generalizes, ECUS head should not habituate to teleportation. Both are pre-registered as v0.14 diagnostics.

\subsection{Phi as the physics of experience -- offline here, online as next step}

There is a temptation to describe experience as prediction error alone. $\Phi$ makes that description precise and then shows why it is incomplete. The same $\epsilon$ can be trivial or telling depending on whether $\Phi$ anticipated it. In the fog, the organism expects to be wrong, so being wrong is not informative. After a teleport in clear conditions, it expected to be right, so being wrong is informative. The ratio $0.13$ operationalizes this. Presence, not error, is the quantity that should drive learning and feeling. This is not a claim about phenomenology. It is a claim about the right engineering target, if one takes predictive processing seriously \cite{friston2010, laukkonen2025}.

We state plainly what $\Phi$ is and is not here. It is a small scalar MSE regressor (15$\rightarrow$64$\rightarrow$1 calibration, 22$\rightarrow$64$\rightarrow$1 integrated) that learned to predict $|\epsilon|$; in Friston's terms it is expected noise, not precision $\Pi$ (inverse conditional variance per channel). What makes it interesting is not its size but its forecast ($r=0.701$, $r_{\mathrm{cross}}=0.730$) and the isolated causal battery ($d=-1.61$). What we do \textit{not} claim is that $\mathrm{presence}=\epsilon/\sigma_\Phi^2$ drives the online 30k loop: it does not; the loop gates on $z$ and $\mathrm{att}$. The online physics-of-precision claim requires a per-channel $\log \sigma^2_{\Phi,c}$ model and direct coupling of $\mathrm{presence}$ to the ECUS drive, versus an attention-only gate (v0.14 4-arm test, Sec.~\ref{sec:limitations}). Only if arm D (presence-coupled, attention frozen) beats arm C (attention-only) is the ``physics of precision'' claim earned.

\subsection{SER versus DECIR: why the mouth matters and why it does not}

We use the Spanish distinction because it is sharper than the English one. SER is what the organism is: ECUS, predictor, $\Phi$, weights, attention. DECIR is what it says. The mouth translates SER into a sentence, but it never selects the next move. Behaviour with and without the mouth is identical. This is not a decorative choice. It is a direct test of a common dismissal: that any ``report'' of an internal state must be a language trick.

We worried this was a philosophical trick, so we tested the more concrete version. The causal test never involves the mouth: conditions A and B differ only in whether $\Phi$ is coupled to action, so the $d=-1.61$ effect cannot be attributed to language. Conversely, the mouth is a frozen translator: given any constructed internal-state prompt it renders it fluently, which shows that DECIR alone is cheap. The substantive claim is not ``the agent says it is uncertain.'' The substantive claim is ``the agent acts as if it were uncertain, in a way that depends on a calibrated forecast of its own error, and that forecast is verifiable against ground truth.'' The mouth only makes that state legible to a human reader.

In classic terms, this deflates the simplest Chinese Room objection for this system. The room (the LLM) is present, but it is not the agent. Inside the room there is a smaller, non-linguistic organism that already knows how to leave the fog. The room translates after the fact.

\subsection{What the failure of C3 means for concept learning -- H\_A}

The fact that habituation generalizes across opposite teleports is not a footnote. It suggests that the predictor's representation of ``a violation'' is not yet a concept but a prototype. In human language we would say the agent learned ``big jump'' rather than ``jump to the northeast by five.'' That is a real limitation if the goal is rich event understanding, and it is an interesting starting point if the goal is developmental. Formally, H\_A is that the trace in $\Delta W$ is habituation of \textit{displacement magnitude}, not of the contingency $P(\Delta p\mid a)$, which is why C3, C4 and persistence are mutually coherent yet incompatible with ``the organism knows what changed.''

We pre-register the Rankin-minimal test of H\_A vs a vector concept (Sec.~\ref{sec:limitations}): same-modality probes S2 $(-2,-2)$, orthogonal S3 $(+2,-2)$, magnitude S4 $(+4,+4)$, type S5 (interoceptive $E$ inversion), true dishabituation ($S_5\rightarrow$ re-$S_1$) and gap recovery. If S2/S3 stay low and dishabituation does not rebound, the claim must remain ``learning without distinguishing''; if they rebound, v0.14 earns stimulus-specific habituation. That experiment is not in this paper and is pre-registered before it is run. The present result is useful because it constrains what we can honestly mean by ``stimulus specificity'' at this scale and matches Rankin et al.\ 2009 rather than contradicting it.

\section{Limitations}
\label{sec:limitations}

We have already stated six limits upfront. We expand them here so that a reader who only reads this section gets the full picture.

\paragraph{1. Toy scale and world (two systems).} The predictor and $\Phi$ are small MLPs. The world was a grid (30-seed battery, $\pm$5) and is now a 20$\times$20 continuous square with hand-coded food/social locations (v0.12 long run, $+$2). Noise is Gaussian and fog is a vertical half-plane. Nothing here resembles the statistics of natural sensing or the architecture of a brain. We do not extrapolate from this scale to any biological claim. The value of the toy is not realism but control: every number can be ablated and every control can be run on a laptop. Table~\ref{tab:main} separates the two systems explicitly since v0.13.

\paragraph{2. No stimulus specificity at fine granularity (C3 renamed).} C3 (stimulus generalization) shows habituation to one large-displacement violation transfers to the opposite vector ($z$ 1.1 vs 0.9, single-seed pilot). We cannot claim stimulus-specific habituation (Rankin characteristic 7), true dishabituation (characteristic 8: $S_1\rightarrow S_2\rightarrow$ re-$S_1$ rebound), or ISI-dependent spontaneous recovery (characteristic 10). At the granularity of ``motor versus interoceptive versus tactile'' the system may discriminate, but at the vector level it does not---we call this ``learning without distinguishing'' (habituation of $\|\Delta p\|$, not of $P(\Delta p\mid a)$).

\paragraph{3. $\Phi$ is not a precision model and its causal effect is confounded.} $\Phi$ is a scalar MSE regressor to $|\epsilon|$ (threshold $r>0.5$), not a per-channel $\log \sigma^2$ (precision) model $\xi=\Pi\epsilon$. Its $r=0.701$ and $r_{\mathrm{cross}}=0.730$ can arise from a zone proxy ($x>14$). The presence ratio $0.13$ is diagnostic offline; presence does not drive the online 30k loop. The isolated $d=-1.61$ (h6\_phi\_causal) lacks attention, while the continuous-organism fog-time drop ($24.7\%$ vs $36.8\%$, $d=-0.43$) is reproduced by a random attention gate (\texttt{62-correcciones-panel-resultados.md}). The gate threshold $\mathrm{att}_{\mathrm{vis}}<0.35$ alone can explain leaving fog. A 4-arm test is required: A ($\Phi$+attention, current), B ($\Phi$ computed/uncoupled), C (attention-only, $\Phi$ frozen), D (presence $\epsilon/\sigma_\Phi^2$ coupled, attention frozen), same seeds $N=30$, food outside fog. If D beats C, the ``physics of precision'' is earned; if A$\approx$C and B$\approx$D, the effect is the attention/zone gate.

\paragraph{4. EWC/ortho are not load-bearing and the trace is likely low-rank.} EWC ($\lambda=5$, diagonal online Fisher) plus $\lambda_{\mathrm{ortho}}=0.01$ was needed to avoid global habituation in v0.11, but for same-task interference it has no effect (recovery $0.48$ for $\lambda\in\{0,0.5,5,50\}$ in \texttt{cuarteto\_habituacion.py}); EWC would need a distinct task B to be a memory dial. The honest interpretation is a low-rank ``ignore large L2'' $\Delta W$, best shown by SVD of $\Delta W$ (how many singular values explain 90\% of the delta?) and by a factored predictor ($f^{\mathrm{pos}}$ and $f^{H}$ with separate losses, surprise per head). Both are pre-registered.

\paragraph{5. The mouth is translation, not recall.} The LLM mouth renders a prompt we construct from measured internal variables. Its fluency should not be mistaken for free report or autobiographical memory. True reportability would require the organism to generate a linguistic description from its own replay without our prompt scaffolding, and to do so in a way that can be wrong in interesting ways. We have not built that.

\paragraph{6. Benchmark is underpowered and removed from Table~\ref{tab:main}.} MiniGrid Empty-8x8 at $N=5$ with vanilla REINFORCE (organism 1.8\% vs ICM 0.6\% vs RND 2.8\% vs random 0.4\%) is a pilot where RND, the strongest baseline, beats the organism. A proper benchmark needs $N=30$, a common PPO backbone, and DoorKey/FourRooms.

\paragraph{7. Other gaps we have not hidden.} Vision/audition were deliberately discarded (``body without eyes''), so we cannot speak to cross-modal contingencies. The 30-seed battery covers motor violations well; interoceptive/tactile violations were tested single-seed only, not in the full blinded battery. Other Rankin criteria (dishabituation, spontaneous recovery with ISI) and isodirectional controls are pre-registered as the v0.14 Rankin-minimal battery (S1--S5, dishabituation, gap, $N=30$, one continuous world, factored predictor, per-channel $\Phi$, SVD). Vision, EWC across tasks, and coupled presence are pending explicitly, not hidden.

\section{Conclusion and Future Work}

We set out to build the smallest thing that could deserve the word ``organism'' in an engineering sense, and to test it as honestly as we could. The organism predicts its own body conditioned on its own action, monitors a homeostatic drive, weights surprise by a learned estimate of its own reliability, acts on that estimate, and consolidates what it learns in its weights. Each of those statements is now tied to a number and a control. The mouth, which is a separate frozen model, only translates.

We close without claiming consciousness. We close on a smaller, more useful note. If awareness has engineering prerequisites, they include at least these: a generative model of the body, a way to know when that model is trustworthy, and a way to act differently when it is not. All three run here for 30,000 steps on a laptop with no discrete GPU.

What comes next, if we continue, is not a bigger MLP. First, the v0.14 Rankin-minimal battery: one continuous world (food outside fog), factored predictor ($f^{\mathrm{pos}}$, $f^{H}$), per-channel $\log \sigma^2_{\Phi}$, 4-arm gate-vs-presence, S1--S5 probes, true dishabituation, ISI gap, SVD of $\Delta W$, $N=30$ (pre-registered in \texttt{63-preregistro-v014-rankin-phi-factorizado.md}). If S2/S3 stay low and dishabituation does not rebound, the paper's claim stays ``learning without distinguishing'' (a coarse but honest model update); if they rebound, v0.14 earns stimulus-specific habituation. Second, a second agent: the present organism can leave the fog alone. What it cannot yet do is learn from another organism that the fog is there, or teach another organism to avoid it, or negotiate a shared policy when food and clarity conflict. That is where culture begins, and where the tetrahedron would need to be extended with a genuinely social dimension: two ECUS loops coupled through actions that are informative for each other, not only for the physics. We have not built that, and we will pre-register it before we do.

Reproduction and falsification are straightforward: run \texttt{framework/organismo\_final.py} and modify the world (move the fog, place food inside it, freeze $\Phi$, shuffle the actions). A minimal system of this size is meant to be inspected and broken, and that is the point. The code and logs that accompany this preprint are archived on Zenodo (v0.12 doi:\,10.5281/zenodo.22191728, concept doi:\,10.5281/zenodo.22191727; v0.13 will receive a new version DOI under the same concept DOI and this line will be updated). A public GitHub release will follow upon publication. The v0.14 pre-registration (Rankin battery + factored predictor + per-channel $\Phi$ + 4-arm gate vs presence + $\Delta W$ SVD) is in \texttt{63-preregistro-v014-rankin-phi-factorizado.md} in the repo.

\section{Reproducibility Statement}

All numbers in Table~\ref{tab:main} are produced by scripts in \texttt{framework/} and logged as JSON in \texttt{results/}. Key scripts: \texttt{organismo\_final.py} (30k integrated run, v0.12 continuous physics with attention, 13$\rightarrow$64$\rightarrow$6, $\Phi$ 22$\rightarrow$64$\rightarrow$1), \texttt{h6\_selfmodel.py} ($\Phi$ calibration, 15$\rightarrow$64$\rightarrow$1, MSE to $|\epsilon|$), \texttt{h6\_phi\_causal.py} (isolated causal gate, no attention), \texttt{rigor\_controles.py} (C1--C4c, seed 7, grid), \texttt{estadistica\_fase2.py} and \texttt{analisis\_fase2.py} (30-seed grid statistics), \texttt{cuarteto\_habituacion.py} (EWC sweep, recovery 0.48), and \texttt{organismo\_final.py} attention training. H5-bis results are in \texttt{results/h5bis.json}; v0.14 pre-registration is in \texttt{63-preregistro-v014-rankin-phi-factorizado.md} (single continuous world, factored heads, per-channel $\Phi$, 4 arms, Rankin S1--S5, dishabituation, gap, SVD). Instructions to reproduce are in \texttt{README.md} and \texttt{requirements.txt}. No discrete GPU is required. On a MacBook Pro (M4 Pro, Apple silicon) with MPS, the full 30-seed battery finishes in under an hour; the single 30k run finishes in a few minutes. We seed with 7 and vary the seed per run for the 30-seed analyses. Bootstrap CIs use 2,000 resamples. We report every control, including the one that fails and is renamed in v0.13 (C3 generalization).

\paragraph{AI assistance disclosure.} An LLM was used for light copy-editing and \LaTeX\ formatting; all results, analyses, tables, and references were verified by the author. No AI-generated content appears in Results or data tables.

\section{Acknowledgements}

We thank the open-science community whose code and reports we build on, especially the authors of the LFM2.5 report, CheckVLA, Gubernaut, and the Levin lab Training Ecosystems framework. LFM2.5-1.2B is used under Liquid AI's LFM Open License v1.0 (free for research and non-commercial use, and for commercial use below USD 10M annual revenue). This work was done entirely on local hardware in San Jos\'e, Costa Rica, on a MacBook Pro (M4 Pro, Apple silicon), with no cloud budget and no discrete GPU. That constraint was not only economic; it kept the system small enough to understand.

\section*{References}
\begin{thebibliography}{25}

\bibitem{chalmers1995} Chalmers, D.~J. Facing up to the problem of consciousness. \textit{Journal of Consciousness Studies}, 2(3):200--219, 1995.

\bibitem{oregan2001} O'Regan, J.~K. and No\"e, A. A sensorimotor account of vision and visual consciousness. \textit{Behavioral and Brain Sciences}, 24(5):939--973, 2001.

\bibitem{thompson2007} Thompson, E. \textit{Mind in Life: Biology, Phenomenology, and the Sciences of Mind}. Harvard University Press, 2007.

\bibitem{friston2010} Friston, K. The free-energy principle: a unified brain theory? \textit{Nature Reviews Neuroscience}, 11(2):127--138, 2010.

\bibitem{tononi2014} Tononi, G., Boly, M., Massimini, M., and Koch, C. Integrated information theory: from consciousness to its physical substrate. \textit{Nature Reviews Neuroscience}, 17(7):450--461, 2016.

\bibitem{baars1988} Baars, B.~J. \textit{A Cognitive Theory of Consciousness}. Cambridge University Press, 1988.

\bibitem{dehaene2011} Dehaene, S. and Changeux, J.-P. Experimental and theoretical approaches to conscious processing. \textit{Neuron}, 70(2):200--227, 2011.

\bibitem{kirkpatrick2017} Kirkpatrick, J. et al. Overcoming catastrophic forgetting in neural networks. \textit{Proceedings of the National Academy of Sciences}, 114(13):3521--3526, 2017.

\bibitem{pathak2017} Pathak, D., Agrawal, P., Efros, A.~A., and Darrell, T. Curiosity-driven exploration by self-supervised prediction. \textit{ICML}, 2017.

\bibitem{burda2018} Burda, Y., Edwards, H., Storkey, A., and Klimov, O. Exploration by random network distillation. \textit{ICLR}, 2019 (arXiv 1810.12894, 2018).

\bibitem{kernelan2014} Keramati, M. and Gutkin, B. Homeostatic reinforcement learning for integrating reward collection and physiological stability. \textit{eLife}, 3:e04811, 2014.

\bibitem{man2016} Man, K. and Damasio, A. Homeostasis and soft robotics in the design of feeling machines. \textit{Nature Machine Intelligence}, 1:446--452, 2019 (preprint 2016).

\bibitem{gillard2022} Gillard, M., Fredes, J., and Rano, I. Intrinsic motivation for habituation in IDSM agents. \textit{IDSM Workshop}, 2022.

\bibitem{levin2026} Samanta, A., Hazan, H., and Levin, M. Training Ecosystems: A Computational Approach to Uncovering Learning Behavior in Unconventional Contexts. \textit{arXiv:2605.30109 [q-bio.PE]}, 2026.

\bibitem{lfm2025} Liquid AI. LFM2 technical report: hybrid SSM-conv at 1.2B scale. \textit{Technical report}, Liquid AI, 2025. Model weights: \texttt{models/LFM2.5-1.2B-MLX-8bit} via \texttt{mlx\_lm}, under the LFM Open License v1.0.

\bibitem{checkvla2026} Liu, Y., Sun, P., Chao, X., et al. CheckVLA: Execution-Time Verification with Action-Conditioned World Model for Long-Horizon Mobile Manipulation. \textit{arXiv:2607.26789 [cs.RO]}, 2026.

\bibitem{gubernaut2026} Sharma, D. Gubernaut: A Deterministic Homeostatic Controller for Affect-Regulated LLM Agents, Validated Across Independent Model Families. \textit{arXiv:2607.24339 [cs.AI]}, 2026.

\bibitem{christovmoore2025} Christov-Moore, L., Juliani, A., Kiefer, A., et al. The Conditions of Physical Embodiment Enable Generalization and Care. \textit{arXiv:2510.07117 [cs.AI]}, 2025.

\bibitem{laukkonen2025} Laukkonen, R., Friston, K., and Chandaria, S. Meditative self-modelling and predictive processing. \textit{Neuroscience \& Biobehavioral Reviews}, 176:106296, 2025.

\bibitem{probe2026} Pavuluri, A., Karkhanis, S., and Mushtaque, U. Asleep at the Wheel: JEPA's Limitations in Evaluating Novel Driving Data. \textit{arXiv:2608.01336 [cs.CV]}, 2026.

\bibitem{thompson1966} Thompson, R.~F. and Spencer, W.~A. Habituation: a model phenomenon for the study of neuronal substrates of behavior. \textit{Psychological Review}, 73(1):16--43, 1966.

\bibitem{rankin2009} Rankin, C.~H., Abrams, T., Barry, R.~J., et al. Habituation revisited: An updated and revised description of the behavioral characteristics of habituation. \textit{Neurobiology of Learning and Memory}, 92(2):135--138, 2009.

\end{thebibliography}

\end{document}
